\subsection{Algoritmo de Shor: Periodicidad y Concentración Espectral}
\label{subsec:patrones_shor}

Shor combina una etapa de exponenciación modular, que introduce estructura periódica en el registro, con una QFT posterior que redistribuye esa información en el dominio de frecuencias. Desde el punto de vista de compresibilidad, esto sugiere dos regímenes distintos dentro de un mismo algoritmo. Antes de la QFT cabe esperar estados con regularidad aritmética y, según la instancia, con subconjuntos más estructurados que el vector completo; después de la QFT esa regularidad puede traducirse en una distribución más densa y con más variación de fase.

En consecuencia, para una estrategia estrictamente sin pérdida no es prudente afirmar un comportamiento uniforme de Shor. Lo más defendible es señalar que algunas etapas podrían resultar relativamente más favorables que otras: las porciones donde la periodicidad induce patrones repetitivos o ceros exactos serían mejores candidatas para compresión, mientras que las etapas donde la QFT densifica el estado tenderían a reducir esa ventaja.
